\section*{الفصل \ref{ch:b2:reproductive-hormones} --- التحكم الهرموني في التكاثر}

\begin{solution}{exo:b2:reproductive-hormones:1}
الوطاء (GnRH، في نبضات) $\to$ النخامى الأمامية
(LH وFSH) $\to$
\omterm{def:b2:mammal-reproduction:testis}{الخصية}: فهرمون LH يؤثر في خلايا لايديغ (التستوستيرون)، وFSH في خلايا سرتولي
(\omterm{def:b2:mammal-reproduction:testis}{تكوّن النطاف}، والإنهيبين). والتغذية الراجعة: التستوستيرون يثبّط
الهرمون المطلق وLH؛ والإنهيبين يثبّط FSH.
\end{solution}

\begin{solution}{exo:b2:reproductive-hormones:2}
المبيض: \omterm{def:b2:reproductive-hormones:cycle}{الطور الجريبي}، الأيام 1--14، والإستراديول من \omterm{def:b2:mammal-reproduction:ovary}{الجريب}
النامي؛ \omterm{def:b2:mammal-reproduction:ovary}{والإباضة}، اليوم 14، \omterm{def:b2:reproductive-hormones:cycle}{واندفاع LH}؛ \omterm{def:b2:reproductive-hormones:cycle}{والطور الأصفري}، الأيام 15--28،
والبروجستيرون من \omterm{def:b2:mammal-reproduction:ovary}{الجسم الأصفر}. والرحم: الطمث، الأيام 1--5
(بسحب الستيرويدات)؛ والطور التكاثري، الأيام 6--14 (الإستراديول)؛
والطور الإفرازي، الأيام 15--28 (البروجستيرون).
\end{solution}

\begin{solution}{exo:b2:reproductive-hormones:3}
يبقي الإستروجين والبروجستين اليوميان LH وFSH منخفضين بالتغذية
الراجعة السالبة، فلا يُجنّد \omterm{def:b2:mammal-reproduction:ovary}{جريب}، ولا يقع ارتفاع إستراديول ولا
ينطلق الاندفاع أبدًا؛ ويثخّن البروجستين المخاط أيضًا. وفي
أسبوع الانقطاع تهبط الستيرويدات وتنسلخ \omterm{def:b2:reproductive-hormones:cycle}{بطانة الرحم} المبنية
تحتها --- وهو نزف سحب، لا طمث بعد
\omterm{def:b2:mammal-reproduction:ovary}{إباضة}.
\end{solution}

\begin{solution}{exo:b2:reproductive-hormones:4}
يرتفع LH وFSH أضعافًا (لانعدام تغذية التستوستيرون أو الإنهيبين الراجعة)؛
ويضمر العضل والرغبة الجنسية والغدد الملحقة. ويعيد التستوستيرون
الصفات ويخفض LH من جديد، لا الخصوبة: فالخصيتان
ذهبتا، وحتى عند رجل له خصيتان يكبت التستوستيرون
الخارجي LH ومن ثم التستوستيرون داخل \omterm{def:b2:mammal-reproduction:testis}{الخصية}
الذي يحتاجه \omterm{def:b2:mammal-reproduction:testis}{تكوّن النطاف}.
\end{solution}

\begin{solution}{exo:b2:reproductive-hormones:5}
$24/1.5 = 16$ نبضة في اليوم؛ و$24/4 = 6$. وحين يموت \omterm{def:b2:mammal-reproduction:ovary}{الجسم الأصفر}
تكون النبضات لا تزال بطيئة، مفضّلةً FSH: فيرتفع FSH أولًا --- وهو
الهرمون الصحيح، إذ إن FSH هو الذي يجنّد المجموعة التالية من
الجريبات.
\end{solution}

\begin{solution}{exo:b2:reproductive-hormones:6}
\qty{200}{pg/mL} $= \qty{200}{ng/L}$: ومنه $200\times 10^{-9}/272 =
\qty{7.4e-10}{mol/L} = \qty{740}{pmol/L}$. و\qty{15}{ng/mL} $=
\qty{15}{\micro g/L}$: ومنه $15\times 10^{-6}/314 = \qty{48}{nmol/L}$. وفي
ميكرولتر واحد: $7.4\times 10^{-16}$ مول $= 4.4\times 10^{8}$
جزيء.
\end{solution}

\begin{solution}{exo:b2:reproductive-hormones:7}
\omterm{def:b2:mammal-reproduction:ovary}{الإباضة} في اليوم $35 - 14 = 21$؛ والخصوبة من اليوم 18 إلى اليوم 22. وفي
الدورات غير المنتظمة لا يمكن التنبؤ \omterm{def:b2:reproductive-hormones:cycle}{بالطور الجريبي}، ومن ثم بيوم
\omterm{def:b2:mammal-reproduction:ovary}{الإباضة}، من الدورات الماضية، فتُفوَّت النافذة أو
تُوضع في غير موضعها.
\end{solution}

\begin{solution}{exo:b2:reproductive-hormones:8}
من اليوم 21 إلى 26 مضاعفتان ونصف: أي $2^{2.5} = 5.7$ وحدة دولية في اللتر --- وهذا كافٍ تمامًا.
\omterm{def:b2:mammal-reproduction:implantation}{والانغراس} في اليوم 25 يعطي $1.4$ وحدة دولية في اللتر في اليوم 26؛ فلا يُنقذ \omterm{def:b2:mammal-reproduction:ovary}{الجسم
الأصفر}، وهو ضامر سلفًا، ويهبط البروجستيرون ويضيع
الجنين مع الدورة. \omterm{def:b2:mammal-reproduction:implantation}{والانغراس} المتأخر سبب شائع
للفقد المبكر.
\end{solution}

\begin{solution}{exo:b2:reproductive-hormones:9}
تستجيب النخامى لنمط الهرمون المطلق لا لمقداره:
فنبضات كل ساعة تصون LH وFSH، والجرعة اليومية نفسها مسرَّبة
باستمرار تلغيهما. وقد استبعدت التجربة علاقة جرعة--استجابة بسيطة.
وشهر من ناهض طويل الأثر: بعد اشتعال أيام قليلة، يهبط LH إلى
ما يقارب الصفر (بتبلّد المستقبلات)، ويهبط الإستراديول إلى
مستويات ما بعد \omterm{prop:b2:reproductive-hormones:overrides}{سن اليأس}، وتضمر \omterm{def:b2:reproductive-hormones:cycle}{بطانة الرحم} ولا تقع دورة --- أي
سن يأس دوائي عكوس.
\end{solution}

\begin{solution}{exo:b2:reproductive-hormones:10}
بلا هرمون: $R^* = 100/0.5 = 200$. ومع الاستمرار: $100/6.5 = 15$. و
نبضة من دقيقتين تزيل $6\times 200\times(2/60) = 40$ مستقبلًا
(أي \qty{20}{\%})؛ وفي الساعة التالية يتقلص العجز بعامل
$e^{-0.5}$، أي تُستعاد منه \qty{39}{\%}. ويستقر المخزون
حيث يساوي الفقد في النبضة الاستعادةَ في الساعة:
أي نحو 150 مستقبلًا قبل كل نبضة، وهو ثلاثة أرباع
الأقصى --- أي عشرة أضعاف المستوى تحت الاستمرار.
\end{solution}

\begin{solution}{exo:b2:reproductive-hormones:11}
\omterm{prop:b2:reproductive-hormones:overrides}{سن اليأس}: لا جريبات، فلا إستراديول ولا إنهيبين، فلا
تغذية راجعة: فيرتفع LH وFSH. وورم البرولاكتين: يبطئ البرولاكتين
مولّد نبضات الهرمون المطلق، فينخفض LH، ولا تُقاد الجريبات،
وينخفض الإستراديول ولا تقع دورة. وورم الخلايا الحُبيبية:
يكبت الإستراديولُ العالي الثابت LH بالتغذية الراجعة السالبة (و
دون \omterm{def:b2:mammal-reproduction:ovary}{جريب} ناضج لا شيء يُباض)، فتتوقف
الدورة مع إستراديول عالٍ --- وتفرط \omterm{def:b2:reproductive-hormones:cycle}{بطانة الرحم} في النمو.
\end{solution}

\begin{solution}{exo:b2:reproductive-hormones:12}
الإستراديول المنخفض أو القصير يخفض الهرمون المطلق وLH؛ أما
الإستراديول العالي الممسوك (فوق \qty{200}{pg/mL} أكثر من 36 ساعة) فيغير
استجابة جملة الكيسبيبتين--الهرمون المطلق واستجابة النخامى،
وقد خزّنت في تلك الأثناء LH، فيثير الهرمون نفسه الآن
دفقة. \omterm{def:b2:mammal-reproduction:ovary}{والإباضة} لا بد أن تكون كلّية أو معدومة وموقوتة: فالبويضة
لا بد أن تُطلق مرة واحدة، في يوم واحد، حين ينضج \omterm{def:b2:mammal-reproduction:ovary}{جريب}. أما المقياس
المتدرج (LH متناسب مع الإستراديول) فكان ليعطي تلوتنًا تدريجيًا جزئيًا،
وجريبات تبيض نصف ناضجة أو عدة جريبات معًا،
وبلا يوم خصيب محدد.
\end{solution}

\begin{solution}{pb:b2:reproductive-hormones:1}
\textbf{1.} الجريبي، الأيام 1--14: الإستراديول يرتفع، والبروجستيرون
منخفض. والأصفري، الأيام 15--28: البروجستيرون عالٍ.
\textbf{2.} بلغ الاندفاع ذروته في اليوم 14 (وبدأ في اليوم 13)؛
\omterm{def:b2:mammal-reproduction:ovary}{والإباضة} بعد نحو 36 ساعة من البداية، في اليوم 15.
\textbf{3.} من اليوم 15 إلى 28: أي 14 يومًا، وهو عمر \omterm{def:b2:mammal-reproduction:ovary}{الجسم
الأصفر}.
\textbf{4.} يرفع البروجستيرون نقطة ضبط تنظيم الحرارة بمقدار
\qty{0.3}{\celsius} إلى \qty{0.4}{\celsius}؛ ويظهر الارتفاع بعد يوم
أو يومين من \omterm{def:b2:mammal-reproduction:ovary}{الإباضة}، فهو يؤكد وقوع \omterm{def:b2:mammal-reproduction:ovary}{الإباضة} لكنه
لا يستطيع الإعلان عنها.
\textbf{5.} الأيام 11 إلى 14 (210، و260، و310، و220): أي ثلاثة أيام إلى أربعة،
أي أكثر من الست والثلاثين ساعة التي يقتضيها المفتاح.
\textbf{6.} الأيام 12 إلى 16.
\textbf{7.} $310\times 10^{-9}/272 = \qty{1.14e-9}{mol/L} =
\qty{1140}{pmol/L}$؛ و$16\times 10^{-6}/314 = \qty{51}{nmol/L}$.
\textbf{8.} $65/8 = 8$. ومع عمر نصف قدره ساعة، ينتصف LH كل
ساعة حين يتوقف الإفراز، فيعود بعد يوم إلى ما يقارب الأساس.
\textbf{9.} \omterm{def:b2:mammal-reproduction:ovary}{الإباضة} نحو اليوم 22، ودورة طولها نحو 36 يومًا.
\textbf{10.} \omterm{def:b2:mammal-reproduction:implantation}{الانغراس} في اليوم 22؛ وبحلول اليوم 27 (أي خمسة أيام، ومضاعفتان
ونصف) نحو \qty{5.7}{IU/L}.
\textbf{11.} هبوط البروجستيرون إلى \qty{2}{ng/mL} في اليوم 27 يعني
أن \omterm{def:b2:mammal-reproduction:ovary}{الجسم الأصفر} يضمر في موعده: فلم ينقذه شيء.
\textbf{12.} البروجستيرون نحو \qty{20}{ng/mL} وهو يرتفع،
والإستراديول نحو \qty{200}{pg/mL}.
\textbf{13.} في نهاية الدورة يهبط الإستراديول والبروجستيرون و
الإنهيبين معًا فيفلت FSH من تغذيته الراجعة؛ ومع نمو المجموعة،
يرتفع الإستراديول والإنهيبين ويهبط FSH، ولا يبقى ناميًا على FSH
الهابط إلا \omterm{def:b2:mammal-reproduction:ovary}{الجريب} الأكثر مستقبلات لهرمون FSH (أي الأكثر خلايا حُبيبية)
--- وتموت البقية.
\textbf{14.} $R^* = 200$ بلا هرمون؛ و$100/6.5 = 15$ تحت
الهرمون المستمر.
\textbf{15.} $6\times 200\times 2/60 = 40$ مستقبلًا، أي \qty{20}{\%}؛
وثابت الزمن $1/d = \qty{2}{h}$، فتُستعاد \qty{39}{\%} من العجز
في ساعة؛ ويستقر المخزون حول 150 مستقبلًا، وهو ثلاثة
أرباع الأقصى.
\textbf{16.} النبضات تبقي النخامى عند نحو 150 مستقبلًا فتجيب
كل نبضة بهرمون LH؛ أما التسريب المستمر فيدفع المخزون إلى
15 فتخرس الخلايا نفسها، مهما كانت الجرعة؛ والنبضات تعيد
المخزون والاستجابة.
\textbf{17.} اليوم الأول: يرتفع LH والتستوستيرون (فالناهض
ينبه مستقبلات لا تزال موجودة --- وهو الاشتعال). وبعد شهر:
تتبلّد المستقبلات، فيقارب LH الصفر، ويصير التستوستيرون عند مستويات
الخصاء، وهذا هو الهدف العلاجي.
\textbf{18.} يرتفع FSH (لانعدام الإنهيبين)، ويبقى LH طبيعيًا؛ فينجو
من الانتقاء جريبات أكثر في كل دورة، مع إباضات متعددة و
توائم ثنائية اللاقحة.
\textbf{19.} في الدورة الطبيعية يقتل هبوط FSH كل الجريبات إلا
المسيطر؛ أما إمداد FSH الثابت فيبقي المجموعة كلها
نامية إلى النضج.
\textbf{20.} تصنع عشرة جريبات أكثر بكثير من \qty{200}{pg/mL} من
الإستراديول مبكرًا، وهو ما كان ليطلق \omterm{def:b2:reproductive-hormones:cycle}{اندفاع LH} سابقًا لأوانه و
إباضةً قبل الجمع؛ والمضادّ يسد الهرمون المطلق فلا يمكن
أن يقع اندفاع.
\textbf{21.} في الجدول بلغ الاندفاع ذروته في اليوم 14 وتبعته \omterm{def:b2:mammal-reproduction:ovary}{الإباضة}
بعد نحو 36 ساعة من بدايته؛ وhCG يحل محل
الاندفاع، والجمع بعد 35 ساعة يلتقط البويضات ناضجة لكنها
لا تزال في \omterm{def:b2:mammal-reproduction:ovary}{الجريب}.
\textbf{22.} $10\times 0.7 = 7$ أجنة، و$7\times 0.4 = 2.8$
\omterm{def:b2:mammal-reproduction:implantation}{كيسة أريمية}؛ ومنه $1 - 0.65^{2} = 0.58$.
\textbf{23.} الإستروجين والبروجستين الثابتان يبقيان FSH وLH منخفضين: فلا
ينمو \omterm{def:b2:mammal-reproduction:ovary}{جريب}، ولا يرتفع الإستراديول إلى العتبة، ولا اندفاع،
ولا \omterm{def:b2:mammal-reproduction:ovary}{جسم أصفر}، ولا ارتفاع بروجستيرون.
\textbf{24.} لا بد من إعطاء التستوستيرون للصفات الثانوية
والرغبة الجنسية؛ لكنه لا يعيد التركيز العالي داخل \omterm{def:b2:mammal-reproduction:testis}{الخصية}
الذي كانت توفره خلايا لايديغ، فيبقى \omterm{def:b2:mammal-reproduction:testis}{تكوّن النطاف}
مكبوتًا --- وهذا هو المقصود.
\textbf{25.} \omterm{def:b2:mammal-reproduction:ovary}{الإباضة} في اليوم 15؛ \omterm{def:b2:reproductive-hormones:cycle}{والطور الأصفري} 14 يومًا؛ والنافذة
الخصيبة الأيام 12--16؛ وأطلق الاندفاعَ إستراديولٌ فوق
\qty{200}{pg/mL} أكثر من 36 ساعة، تحقق في الأيام 11--14.
\end{solution}
